% Atomic-response to finite-volume Bill transfer and Coffer price conversion.
\section{Atomic responses and the finite-volume root price}
\label{sec:ex-lattice-trace-coffer-supplier}

This section records the actual transfer used by the public trace supplier.
There are two distinct operations. An atomic response first becomes a
finite-volume response with a root-price estimate. The price estimate is
then converted into the two components of a Bill. We give the conversions
and the quantitative endpoint of the Coffer argument, while keeping the
upstream response, paid-price, and price-duality receipts explicit.

\subsection{From an atomic response to the Coffer input}

Let \(X\) be an order-zero coefficient at scale \(k\), \(g\) its local field,
and \(Z=\widehat\beta_k\). Write \(x_X\) for its coordinate, \(X^{\rm f}\) for its forced part, and \(D_X\) for its density. The atomic certificate supplies a nonnegative
remainder allowance \(s\), a response \(L\) with attainment cost at most
\(e\), and
the identities
\begin{align}
 \|x_X\|_{\rm coeff}+s&\leq q, \label{eq:ex-trace-atomd-charge}\\
 \operatorname{profile}(X^{\rm f})&=\operatorname{hook}(L),\\
 D_X(U)&=\operatorname{eval}(U)\,g(U).
 \label{eq:ex-trace-atomd-data}
\end{align}
The remainder bound is measured at scale \(k\), at order zero and with
tolerance \(\sqrt{\operatorname{forcedAllowance}_k}\). The statement that
the response \(L=(L_Y)_Y\), indexed by supports \(Y\), has attainment cost
at most \(e\) is a certificate, not a scalar equality: every component
\(L_Y\) is measurable and gauge invariant, vanishes at every configuration
that is healthy on \(Y\) at the threshold of the scale-\(k\) budget, and
takes the same value at two configurations whose links on \(Y\) agree and
whose threshold flags in \(Y\) agree; \(L_Y\neq0\) only if \(Y\) contains
the marked origin, and there are numbers \(a_Y\geq0\) with
\[
 |L_Y(U)|\leq a_Y\ \text{ for all }U,\qquad
 \sum_Y e^{(\alpha_k+2)\operatorname{cov}_k(Y)}a_Y\leq e,
\]
where \(\alpha_k\) is the support exponent and \(\operatorname{cov}_k\)
the cover number of that budget. Set
\(u=\operatorname{alto}_{k,d}\), the inverse field dilation, and \(G'=ug\).
Multiplying the density identity by \(u\) gives
\[
 uD_X(U)=\operatorname{eval}(U)G'(U).
\]
Thus the atomic response becomes the Coffer input, with unchanged format,
remainder and forced-profile data.

The selected atomic charge budget satisfies
\begin{equation}
 4(1+c_k)q\leq Z^3,
 \label{eq:ex-trace-coffer-budget}
\end{equation}
where \(c_k\geq1\) is the selected column allowance of the history at
scale \(k\), a real parameter of the stage data and not a cardinality.
Hence
\[
 4(1+c_k)(\|x_X\|_{\rm coeff}+s)\leq Z^3.
\]
The second needed comparison is \(Z\operatorname{fly}_d(k)\leq Z^2\),
using \(\operatorname{fly}_d(k)\leq Z\) and \(Z\geq1\). The price construction
chooses the atomic charge allowance so both comparisons hold.

The endpoint price-duality estimate used below has the following exact
meaning. A response duet \(P\) runs from a level \(\ell\) to \(k\leq N\)
in the same selected history. At each intermediate level its coefficient
has admissible order-zero format and remainder bound; its marked site is the
spatial origin and these sites are related by the one-step scale map. The
coefficient and cost changes satisfy the paid direction inequality, and
the densities satisfy the exact one-step transfer identity. Its left
endpoint coefficient vector is zero, and its terminal cost obeys the
spread estimate
\[
 \operatorname{cost}_k\leq
 w+c_k\bigl(w+\|P_k.x\|_{\rm coeff}\bigr),\qquad w\geq0.
\]
Write \(\kappa_k>0\) for the square root of the level-\(k\) monitor.
The endpoint estimate, which is required for all sufficiently large
period indices, uniformly in the volume index
\(N\geq k\), in the paid comparison and in \(w\geq0\), is
\begin{equation}
 \kappa_k\,\|P_k.x\|_{\rm coeff}
 \leq4(1+c_k)w.
 \label{eq:ex-trace-sharp-read}
\end{equation}
The zero in the left-end condition is a coefficient vector; the marked sites
are at the spatial origin, not zero-valued observations. The finite recurrence
and adjacent-level estimate are given in
Section~\ref{sec:ex-lattice-trace-price-recurrence}. The paid-cone estimate is supplied by the canonical argument in
Section~\ref{sec:reference-cone}, under its explicit direction, room, bank,
and endpoint-spread assumptions. None of these bounds follows from the atomic
response alone.

\subsection{The finite Coffer transfer}

We keep the source parameters distinct. As above, \(c_k\geq1\) is the
selected column allowance at level \(k\). Let \(X\) be the input coefficient, \(S\) its remainder budget,
\(E\) the paid response and \(e\) its attainment cost, and
\(\zeta_i>0\) its payment allowance. Let \(a\in(0,1]\) be the mask fraction,
\(v'\geq0\) the forced-allowance budget, and \(q,B>0\) two coefficient budgets.
The actual Coffer hypotheses include
\begin{equation}
 \|x_X\|_{\rm coeff}+S\leq q,\qquad
 4(1+c_k)q\leq B,\qquad e\zeta_k\leq v'B.
 \label{eq:ex-trace-coffer-inputs}
\end{equation}
They also include exact coefficient format, remainder, payment, density and
forced-profile identities; scale growth and one-step response receipts; and
the melting condition. If \(B_k\) is the spatial-spacing parameter of the
scale budget and \(R_k=\operatorname{siteFactor}(B_k)\), the latter requires
support exponent at least four and
\[
 t=\operatorname{forcedAllowance}_k
  +R_k\sqrt{\operatorname{forcedAllowance}_k}+a+R_kv',\qquad
 (2B_k+1)^4t\leq3,\qquad e^t-1\leq1.
\]
Here \(B_k\) is not the coefficient budget \(B\) in
\eqref{eq:ex-trace-coffer-inputs}.

For \(J\leq m<k<N\), let \(\operatorname{gain}(k,m)>0\) denote the
scale gain from \(m\) to \(k\). The output weight is
\begin{equation}
 w_m=a\,\frac{\zeta_m/2}
 {\|\operatorname{alto}_{m,d}\|(B/\operatorname{gain}(k,m))}.
 \label{eq:ex-trace-coffer-weight}
\end{equation}

For clarity, if \(\mathcal T_J\) is the set of reflected scale-\(J\) cells
in the volume-\(N\) tiling and \(K=|\mathcal T_J|\), with \(\mu_{N,p}\) the normalized Wilson measure, the root functional is
\[
 \operatorname{root}_{J,N,p}(F)=
 \left\|\int\prod_{x\in\mathcal T_J}
       \operatorname{copy}_x(F)(U)\,d\mu_{N,p}(U)\right\|^{1/K}.
\]
Here \(\operatorname{copy}_x(F)\) evaluates \(F\) on the translated
reflection half of the configuration; when that reflection reverses
orientation it also complex-conjugates the value. This is the
reflection-positive tiling functional, with its \(K\)th-root normalization.

Assume \(N>k\) and that the reference level is no larger than \(k\).
The response data are an admissible
sequence of coefficient voices from the reference level through \(N\), an
admissible top coefficient, an exact paid response identity between each
adjacent pair, and the terminal forced-profile identity. The input field
has the format, remainder, attainment, and density identities stated above.
Assume also the melting and payment inequalities and the endpoint estimate
\eqref{eq:ex-trace-sharp-read}. Put
\[
 \widetilde G'=\operatorname{back}_{k\to N}G'.
\]
Under these hypotheses the finite construction returns one public label
\(b\), chosen before any window, with the following two properties. First,
\begin{equation}
 b'\,\|x_X-\mathcal C_k(b)\|_{\rm coeff}
   \leq 4(1+c_k)q. \label{eq:ex-trace-coffer-column}
\end{equation}
Here the window from scale \(J\) to scale \(m\) is called \emph{open} if
\[
 J\leq m,\qquad \sigma_m^2\leq A_J/A_m,
\]
where \(\sigma_m\) is the cover size of the selected clock at level \(m\)
and \(A_j\) the field dilation. Second, for the same \(b\) and for every
triple \((J,m,l)\) such that the
scale-\(J\)-to-\(m\) window is open, the reference level is no larger than
\(m\), \(m<k\), the scale \(J\) has room for the radius \(l\) at scale
\(k\), and \(G'\) is interior to the window of radius \(l\),
\begin{equation}
 \operatorname{root}_{J,N,p}(\widetilde G'-\mathcal O_{N,p}b)
   \leq \frac{\operatorname{stakes}(J,m)}{w_m}.
 \label{eq:ex-trace-coffer-root}
\end{equation}
The order of quantifiers matters: a label depending on the window would
not give a Bill, whose root estimate is required for all windows with one
label. Here \(b'\) is the endpoint price-duality constant and
\(\operatorname{root}_{J,N,p}\) is the reflection-positive root functional
on the volume-\(N\) field, defined using the scale-\(J\) tiling. In this
application the endpoint constant $b'$ is $\kappa_k$. The column
\(\mathcal C_k(b)\) is defined through the bare reference coordinate of
the public label: if \(\mathsf T\) denotes the terminal identification of
the selected response hierarchy, a linear isomorphism of the grade-\(d\)
coefficient space onto itself which sends bare reference
coordinates \(v\) to public labels, and \(I_k v\) is the coefficient of
the calibrated primary field \(\mathcal P_k(v)\) at level \(k\), then
\[
 \mathcal C_k(b)=I_k(\mathsf T^{-1}b).
\]
The construction below produces a bare coordinate \(v\) and sets
\(b=\mathsf T v\); the column estimate is proved for \(I_kv\), while the
observation \(\mathcal O_{N,p}b\) in \eqref{eq:ex-trace-coffer-root} is
that of the public label. That the level-\(k\) entries of the selected
response hierarchy, read at \(\mathsf T^{-1}b\), are these calibrated
primary fields is an identity of the selected hierarchy, used here as part
of the response data.

For the column estimate, apply the signed reference-difference construction
of Section~\ref{sec:reference-difference} to the Coffer coefficient field
\(X\) at terminal level \(k\), with its remainder bound \(S\) and total
charge \(q\). The construction assumes one fixed background history and
calibrated primary fields \(\mathcal P_i(v)\) whose coefficient is
\(I_i v\), together with the bank and column inequalities
\eqref{eq:reference-bank} and the exact degree-\(d\) density transition.
At each backward step it first transports the current zero-coefficient
remainder by the zero-coordinate transition
\eqref{eq:reference-zero-transition}; it then subtracts the primary field
with the resulting coefficient, so that the new remainder again has
coefficient zero. Linearity cancels the two primary densities exactly, and
the zero-coordinate transition supplies the density of the transported
remainder. The paid direction and remainder estimates
\eqref{eq:reference-direction} apply to the resulting signed difference,
not to an unsigned estimate obtained by subtracting costs.

The top family has terminal charge at most \(q\), since
\(\|x_X\|_{\rm coeff}+S\le q\). Choose its single bare reference coordinate \(v\) by
setting the left endpoint coefficient to zero, and put \(b=\mathsf Tv\).
The output is precisely
\eqref{eq:reference-produced-difference}: for the chosen coordinate, the
terminal coefficient is \(x_X-I_kv=x_X-\mathcal C_k(b)\), the terminal
density is
\(D_X-D_{\mathcal P_k(v)}\), and the left endpoint coefficient vector
vanishes. This choice does not involve any window \((J,m,l)\). Every intermediate field and the selected reference field are
evaluated against the same finest volume and background history; this is
what permits the exact density cancellation. At the terminal level the reference construction also retains its spread
bound \(r_k\le q+c_k(q+\|Z_k.x\|_{\rm coeff})\), the precise spread
condition in the Woke duet. Thus it applies \(\operatorname{SharpRead}\)
with budget \(w=q\). The stronger full terminal estimate in
\eqref{eq:reference-produced-difference}, with \(C_k=c_k\), gives
\(r_k+\|Z_k.x\|_{\rm coeff}\le4(1+c_k)q\); this is the coarse
transport budget used by the subsequent construction. The signed dualization
then gives the zero-site paid-response duet required by SharpRead, and its
endpoint inequality is exactly \eqref{eq:ex-trace-coffer-column}. Here zero refers to the left coefficient
vector and to the duet's marked spatial site being the origin. It does not
assert that the observed field values vanish.
For the rooted estimate, fix a window \((J,m,l)\) as above; the label
\(b\) is already fixed. The same dual chain is planted once at each site of
the finite scale-\(J\) census. The construction couples each translated
negative endpoint to the local lens response, then propagates the resulting
polynomial to level \(m\). The properties of this planted expansion that
are used below are an input from the response construction, not proved in
this section; we state them explicitly. \emph{Planted-expansion input.}
Let \(V\) be the set of census sites, \(K=|V|\), with marks \(o_a\),
\(a\in V\), in the level-\(m\) lattice. The level-\(m\) census polynomial
\(\mathsf X(U)\) is integrable against the reference
measure and has the following structure. There are singleton site
amplitudes \(\Phi_a(U)\) and singleton box amplitudes \(\omega_{a,P}(U)\),
\(a\in V\), \(P\) a box, with
\[
 w_a|\Phi_a(U)|\leq2,\qquad \sum_P w_a|\omega_{a,P}(U)|\leq1
 \qquad\text{for every }a\text{ and every }U,
\]
and two correction families, an ordinary family \(H_Y(U)\) indexed by
supports \(Y\) and a flagged family \(K_P(U)\) indexed by boxes \(P\), such
that
\[
 \mathsf X(U)=e^{E_m(U)}\,
   \operatorname{pass}\Bigl(\sum_a\Phi_a(U)\,\mathsf t_a+\sum_YH_Y(U)\Bigr)
   \cdot\operatorname{field}\Bigl(P\mapsto
       \mathsf b_P(U)+\sum_a\omega_{a,P}(U)\,\mathsf t_a+K_P(U)\Bigr),
\]
where \(E_m\) is the exponent of the level-\(m\) effective density,
\(\mathsf t_a\) is the census monomial of the site \(a\), \(\mathsf b_P\)
is the box activity term of the level-\(m\) representation, and
\(\operatorname{field}\) is evaluated on the flagged set of \(U\). Each
correction family satisfies the anchored support-cost bound with constant
\(1\) at the support exponent \(\alpha_m\) of the level-\(m\) budget: there
are numbers \(h_Y\geq0\) bounding the weighted absolute coefficient sums,
\(\mathcal P_w(H_Y(U))\leq h_Y\) for all \(U\), with
\(\sum_{Y\ni x}h_Ye^{\alpha_m\operatorname{cov}_m(Y)}\leq1\) for every
site \(x\), and likewise for \(K_P\) with the support of the box \(P\).
Every nonzero correction term has a census mark \(o_a\) in its support.
Consequently, summing the anchored bound over the \(K\) marks,
\begin{equation}
 \sum_Y\mathcal P_w(H_Y(U))\leq K,\qquad
 \sum_P\mathcal P_w(K_P(U))\leq K
 \qquad\text{for every }U.
 \label{eq:ex-trace-coffer-corrections}
\end{equation}
Finally, the
planted coefficient is
the moment in the root functional: integrating it against the level-\(m\)
reference measure gives the volume-\(N\) Wilson partition function times
the normalized moment \(M\) of
\(\prod_{x\in\mathcal T_J}\operatorname{copy}_x(\widetilde G'-
\mathcal O_{N,p}b)\). In the source these properties are derived from the
same response hierarchy and its adjacent paid-response identities, the
terminal forced-profile identity, the melting condition, occupancy and
placement conditions on the planted sites at the levels \(m\) and \(k\),
the two payment inequalities, and integrability of the initial and coupled
polynomials; the moment identity follows at each transfer from
normalization of the transition kernel and the exact response-density
identity, and at the terminal level from the definition of the marked
coefficient.

The polynomial-price calculation behind the finite-gas estimate is as
follows. For a finite census polynomial \(Q\), write
\[
 \mathcal P_w(Q)=\sum_s\operatorname{toll}(w,s)\|Q_s\|
\]
for its weighted absolute coefficient sum. The price is submultiplicative.
For box terms \(K_p\), the finite sum over disjoint-support patterns obeys
\[
 \mathcal P_w(\operatorname{field}(K))
 \leq\prod_p(1+\mathcal P_w(K_p))
 \leq\exp\!\left(\sum_p\mathcal P_w(K_p)\right),
\]
because each pattern's product price is at most the product of its factors
and \(1+x\leq e^x\). The same submultiplicativity, applied term by term to
the finite truncated exponential, gives
\(\mathcal P_w(\operatorname{pass}(Q))\leq e^{\mathcal P_w(Q)}\).
The singleton site and box amplitudes contribute \(2K\) and \(K\) to the
exponent, by the two singleton bounds of the planted-expansion input. The
ordinary and the flagged correction families contribute \(K\) each, by
\eqref{eq:ex-trace-coffer-corrections}. The total is \(5K\); the box
activity terms \(\mathsf b_P\) and the exponent \(E_m\) are paid by the
finite-volume denominator bound used next.

By the planted-expansion input, the integral of the selected coefficient
equals the
volume-\(N\) Wilson partition function times the normalized moment \(M\).
The partition function is positive. Dividing by it and using the finite
volume denominator bound yields
\begin{equation}
 |M|\,\operatorname{toll}(w,\mathfrak s)
 \leq \exp\!\left(|\Lambda_m|\operatorname{Treasury}(m)
                         +K(2+1+2)\right).
 \label{eq:ex-trace-spiker}
\end{equation}
This is an absolute-value estimate and does not use cancellation. Its
hypotheses include the response construction's support-cost and moment
identity receipts.

In the Coffer application all site weights equal \(w_m\), so
\(\operatorname{toll}(w,\mathfrak s)=w_m^K\). The tiling identity is
\[
 |\Lambda_m|=K\left(\frac{A_J}{A_m}\right)^4.
\]
Taking the positive \(K\)-th root of \eqref{eq:ex-trace-spiker} yields
\[
 |M|^{1/K}\leq
 \frac{\exp((A_J/A_m)^4\operatorname{Treasury}(m)+5)}{w_m}
 =\frac{\operatorname{stakes}(J,m)}{w_m},
\]
which is \eqref{eq:ex-trace-coffer-root}. The moment identity is needed to
identify this \(M\) with the desired rooted observable.

\subsection{Using the output in the Bill}

The Bill is evaluated on the inflated field \(G'=u g\). Its root error is
\(\widetilde G'-\mathcal O_{N,p}b\), where
\(\widetilde G'=\operatorname{back}_{k\to N}G'\). In the Bill specialization, \(a=Z^{-1}\), \(v'=Z^{-1}\), \(B=Z^3\), and \(e=Z\), while \(\zeta_i=\operatorname{fly}_d(i)\) at each scale. Thus the payment condition is \(e\zeta_k=Z\operatorname{fly}_d(k)\leq Z^2=v'B\). The gas weight
from \eqref{eq:ex-trace-coffer-weight} is
\[
 w_{\rm gas}=Z^{-1}
 \frac{\operatorname{fly}_d(m)/2}
 {\|\operatorname{alto}_{m,d}\|(Z^3/\operatorname{gain}(k,m))}.
\]
The factor $Z^{-1}$ is the mask fraction $a$ in the Coffer input, so this is
exactly the native Bill denominator.

The other two entries of the Bill use the specific atomic charge. Fix
\(u>0\) and let the charge in \eqref{eq:ex-trace-atomd-charge} be
\[
 q=q_k=u\,\vartheta_k^2\operatorname{Age}_k(1/6),
 \qquad
 \operatorname{Age}_k(a)=\exp\bigl(a(\log\log\widehat\beta_k)^2\bigr),
\]
where \(\vartheta_k\) is the threshold of the scale-\(k\) budget; thus
\(\kappa_k=\operatorname{Age}_k(1/2)\) and
\(\operatorname{Age}_k(a)\operatorname{Age}_k(a')=\operatorname{Age}_k(a+a')\).
The selected schedule has \(c_k\leq\log\widehat\beta_k\). Since
\(\vartheta_k^2\operatorname{Age}_k(1/6)=\widehat\beta_k^{-1}
\operatorname{Age}_k(2+1/6)\) tends to zero faster than any power of
\(\log\widehat\beta_k\) grows, and
\(\log\widehat\beta_k=\exp(\log\log\widehat\beta_k)\) is eventually
smaller than any positive multiple of \(\operatorname{Age}_k(1/12)\), and
\(\widehat\beta_k\to\infty\), for all large \(k\)
\[
 q_k\leq1,\qquad 4(1+\log\widehat\beta_k)q_k\leq Z^3,\qquad
 4(1+\log\widehat\beta_k)u\leq\operatorname{Age}_k(1/12);
\]
the second inequality gives \eqref{eq:ex-trace-coffer-budget}. The norm
entry of the Bill is \(\|x_X\|_{\rm coeff}\leq1\); it follows from
\eqref{eq:ex-trace-atomd-charge}, nonnegativity of the remainder allowance
and \(q_k\leq1\):
\(\|x_X\|_{\rm coeff}\leq\|x_X\|_{\rm coeff}+s\leq q_k\leq1\). The bias
entry is obtained from \eqref{eq:ex-trace-coffer-column} with
\(b'=\kappa_k\) alone:
\[
 \|x_X-\mathcal C_k(b)\|_{\rm coeff}
 \leq\frac{4(1+c_k)q_k}{\kappa_k}
 \leq\operatorname{Age}_k(1/12)\,\vartheta_k^2
      \operatorname{Age}_k(1/6)\operatorname{Age}_k(-1/2)
 =\vartheta_k^2\operatorname{Age}_k(-1/4),
\]
because \(1/12+1/6-1/2=-1/4\). This is the bias bound which
Section~\ref{sec:ex-lattice-trace-bill-prices} takes as a hypothesis on a
paid label; that section does not prove it.

In the source, the resulting statement has the following order of
choices, under the standing hypotheses on the selected family, which
include the endpoint estimate \eqref{eq:ex-trace-sharp-read} with
\(b'=\kappa_k\) as a hypothesis. For every fixed \(u>0\) and grade
\(d\): for all sufficiently
large \(k\), for all sufficiently large periods \(p\) (depending on \(k\)),
for every volume index \(N>k\) and every atomic certificate \((X,g)\) at
scale \(k\) with charge \(q_k\) and attainment cost at most
\(Z=\widehat\beta_k\), there is a public label \(b\) such that
\((g,x_X,b)\) is a Bill, with the root estimate for all windows. The
period threshold is the one at which the endpoint estimate
\eqref{eq:ex-trace-sharp-read} holds with \(\ell\) the reference level.

The outstanding suppliers are the atomic response, the endpoint
price-duality recurrence, and the local action and occupancy estimates
which construct the site/box expansion and its absolute-value envelope.
The argument above spells out the numerical transfer once those receipts
are supplied; it does not claim to have proved those producers.
